\def\baselinestretch{1}

\chapter{Model}\label{ch:model}

\def\baselinestretch{1.66}

\section{Introduction}\label{sec:model_introduction}

% --- Spring/damper pics --------------------------------------------------
% Long-form damper: cylinder body + protruding piston rod, total length 2L.
% Cylinder anchors at left, piston rod exits to the right.
%   #1 = full length (cm).  Piston-rod length is half of total.
\tikzset{
  vdamper/.pic={
    \draw[line width=0.5pt, mnnwInk] (0,-0.45)            -- (0,-0.18);
    \draw[line width=0.5pt, mnnwInk] (-0.11,-0.18) rectangle (0.11,0.18);
    \draw[line width=0.5pt, mnnwInk] (-0.11,0)            -- (0.11,0);
    \draw[line width=0.5pt, mnnwInk] (0,0.18)             -- (0,0.45);
  },
}
% Long horizontal damper between (xa, y) and (xb, y).
% Cylinder is the LEFT 60%, piston rod the RIGHT 40%.
\newcommand{\hdamperLong}[3]{%
  \pgfmathsetmacro{\xa}{#1}\pgfmathsetmacro{\xb}{#2}\pgfmathsetmacro{\yy}{#3}
  \pgfmathsetmacro{\Lcyl}{(\xb-\xa)*0.55}
  \pgfmathsetmacro{\xcylR}{\xa+\Lcyl}
  \pgfmathsetmacro{\xpisR}{\xb}
  % cylinder body (open on the right end)
  \draw[line width=0.5pt, mnnwInk] (\xa,\yy-0.13) -- (\xa,\yy+0.13);% closed left wall
  \draw[line width=0.5pt, mnnwInk] (\xa,\yy+0.13) -- (\xcylR,\yy+0.13);
  \draw[line width=0.5pt, mnnwInk] (\xa,\yy-0.13) -- (\xcylR,\yy-0.13);
  % piston disc inside cylinder (about 70% along)
  \pgfmathsetmacro{\xpiston}{\xa+\Lcyl*0.70}
  \draw[line width=0.5pt, mnnwInk] (\xpiston,\yy-0.11) -- (\xpiston,\yy+0.11);
  % piston rod from disc to right anchor
  \draw[line width=0.5pt, mnnwInk] (\xpiston,\yy) -- (\xpisR,\yy);
}
% Long vertical damper between (x, ya) and (x, yb).
% Cylinder is the BOTTOM 60%, piston rod the TOP 40%.
\newcommand{\vdamperLong}[3]{%
  \pgfmathsetmacro{\xx}{#1}\pgfmathsetmacro{\ya}{#2}\pgfmathsetmacro{\yb}{#3}
  \pgfmathsetmacro{\Lcyl}{(\yb-\ya)*0.55}
  \pgfmathsetmacro{\ycylT}{\ya+\Lcyl}
  \draw[line width=0.5pt, mnnwInk] (\xx-0.13,\ya) -- (\xx+0.13,\ya);
  \draw[line width=0.5pt, mnnwInk] (\xx-0.13,\ya) -- (\xx-0.13,\ycylT);
  \draw[line width=0.5pt, mnnwInk] (\xx+0.13,\ya) -- (\xx+0.13,\ycylT);
  \pgfmathsetmacro{\ypiston}{\ya+\Lcyl*0.70}
  \draw[line width=0.5pt, mnnwInk] (\xx-0.11,\ypiston) -- (\xx+0.11,\ypiston);
  \draw[line width=0.5pt, mnnwInk] (\xx,\ypiston) -- (\xx,\yb);
}

As we established before, the accurate relative positioning between the reticle and the wafer is crucial for \textit{step-and-scan} machines to produce functional ICs. To analyze these positioning requirements in a multi-body setting, we leverage existing literature~\cite{Al-Rawashdeh2022_Caracterization} that proposes modeling the lithography scanner as a combination of three open-loop serial chains: the reticle chain ($j=1$, blue), the optics chain ($j=2$, purple), and the wafer chain ($j=3$, red); the latter being the crux of the problem studied here.

The base frame serves as the common structural root for all three chains. In the reticle chain, a dual-stage architecture is employed to satisfy the competing requirements for travel range and precision. Initially, the long-stroke motor (LS) ---of the planar moving-coil type\cite{tissingh2023improving}--- consisting of four coil actuators arranged in a platform shape facilitates large-range travel (of at least 300\,mm) with moderate micrometer accuracy in 6-DOF. However, being directly mounted to the base frame, it is inherently susceptible to vibrations propagated from the rest of the machine.

To achieve nanometric tracking accuracy, a short-stroke motor (SS) is appended to the long-stroke carrier. Since travel in this stage is limited up to 1\,mm, it is realized with Lorentz actuators which provide contact-free, high-bandwidth positioning. Given that the Lorentz force depends entirely on the current within the actuator and is independent of the stator's motion, it effectively decouples the end-effector from the base frame disturbances through magnetic suspension\cite{Butler2011_PositionControl}. To minimize linearity errors and optimize force transfer, the long-stroke stage carrying the stator must maintain the optimal clearance to the magnets of the rest of the short-stroke actuator. This arrangement is replicated for the wafer chain.

The optics chain is attached to the base frame through a metrology frame which is in turn supported by pneumatic vibration isolators\cite{Butler2011_PositionControl}. While the optics box includes adaptive optics with multiple elements, it is reduced in this framework to a single fixed lens. This allows the model to focus on machine compliance and structural vibration transfer rather than involving active optical control\cite{Al-Rawashdeh2022_Caracterization}.

The physical connections between these stages, which constitute the primary transmission path for vibration, are represented by lumped-parameter spring-damper systems, according to the Kelvin-Voigt model. As its proponents state, this low-fidelity modeling choice is sufficient to study the dynamics between the multiple bodies during \textit{step-and-scan} motion\cite{Al-Rawashdeh2022_Caracterization}.

Figure~\ref{fig:litho_model_full} illustrates the structural topology and hierarchical arrangement within each chain. Horizontal Kelvin–Voigt elements $(K_{i,j},C_{i,j})$ couple each stage laterally to the fixed structural reference depicted as a solid line in the color of the corresponding chain; in particular, elements $(K_1,C_1)$ provide passive vibration isolation between $\mathcal{F}_0\!\rightarrow\!\mathcal{F}_1$ through the solid green line that extends the base frame ($i=1$) into the coupling to the first body ($i=2$) in any $j$ chain. The in-plane deviations $f_{i,j}$ represented by the horizontal blue double-headed arrows measure the relative displacement between each body. The dashed vertical centerline represents the nominal reference for $f^0_{i,j}$, when all relative coordinates are zero.

\begin{figure}[htbp]
\centering
\begin{tikzpicture}[scale=0.5, transform shape,
  every node/.style    ={font=\small},
  rfrm/.style          ={draw=frameEdge, line width=0.6pt, fill=frameFill,
                         minimum height=0.80cm, font=\small\bfseries,
                         inner sep=2pt},
  rblk/.style          ={draw=reticleEdge, line width=0.5pt, fill=reticleFill,
                         minimum height=0.55cm, font=\footnotesize,
                         text=mnnwInk, inner sep=2.5pt},
  rflx/.style          ={draw=reticleEdge, line width=0.4pt, fill=reticleFlex,
                         minimum height=0.42cm, font=\scriptsize,
                         text=mnnwInk, inner sep=2pt, dashed},
  wblk/.style          ={draw=waferEdge, line width=0.5pt, fill=waferFill,
                         minimum height=0.55cm, font=\footnotesize,
                         text=mnnwInk, inner sep=2.5pt},
  wflx/.style          ={draw=waferEdge, line width=0.4pt, fill=waferFlex,
                         minimum height=0.42cm, font=\scriptsize,
                         text=mnnwInk, inner sep=2pt, dashed},
  oblk/.style          ={draw=opticsEdge, line width=0.5pt, fill=opticsFill,
                         minimum height=0.55cm, font=\footnotesize,
                         text=mnnwInk, inner sep=2.5pt},
  oflx/.style          ={draw=opticsEdge, line width=0.4pt, fill=opticsFlex,
                         minimum height=0.42cm, font=\scriptsize,
                         text=mnnwInk, inner sep=2pt, dashed},
  ghost/.style         ={draw=mnnwHair, dashed, line width=0.35pt,
                         fill=none, minimum height=0.55cm},
  hsp/.style           ={line width=0.5pt, mnnwInk, decorate,
                         decoration={zigzag, segment length=3.6pt, amplitude=2pt,
                                     pre length=2pt, post length=2pt}},
  vsp/.style           ={line width=0.5pt, mnnwInk, decorate,
                         decoration={zigzag, segment length=3.6pt, amplitude=2pt,
                                     pre length=2pt, post length=2pt}},
  darr/.style          ={Stealth-Stealth, line width=0.5pt, mnnwSignal},
  lbl/.style           ={font=\scriptsize, text=mnnwInk},
  ftn/.style           ={font=\scriptsize, text=mnnwSignal},
  axis/.style          ={dashed, line width=0.4pt, mnnwMute},
  pillar/.style        ={fill=black!8, draw=mnnwHair, line width=0.3pt},
  hatchfloor/.style    ={pattern=north east lines, pattern color=mnnwMute},
]

% =====================================================================
% Helper macro for one chain body row.
%   #1 y-coord            #2 block-style (rblk/wblk/oblk/rflx/wflx/oflx)
%   #3 previous_x         #4 horizontal offset from Pd (positive = right)
%   #5 block label        #6 K-label   #7 C-label   #8 f-label
% =====================================================================
% Trackers for drawing dashed lines between row couplings
\def\lastSpringY{-100}
\def\lastX{-100}

\def\getfirstchar#1#2\relax{#1}

\newcommand{\chainrow}[8]{%
  % Determine chain color based on style prefix
  \edef\styleName{#2}%
  \def\currColor{mnnwMute}%
  \edef\firstChar{\expandafter\getfirstchar\styleName\relax}%
  \if\firstChar w\def\currColor{waferEdge}\fi
  \if\firstChar o\def\currColor{opticsEdge}\fi
  \if\firstChar r\def\currColor{reticleEdge}\fi

  % Draw dashed line between previous row's spring and current row's damper at the anchor
  \pgfmathparse{abs(#3 - \lastX) < 0.001 ? 1 : 0}
  \ifnum\pgfmathresult=1
    % Use frameEdge if at the centerline (connection to base), else use chain color
    \pgfmathparse{abs(#3) < 0.001 ? 1 : 0}
    \ifnum\pgfmathresult=1
      \def\lineColor{frameEdge}
    \else
      \def\lineColor{\currColor}
    \fi
    \draw[line width=0.35pt, \lineColor] (#3, \lastSpringY) -- (#3, #1-0.62);
  \fi

  % 1. Render the real block first to determine its natural width
  % We use 'inner xsep=6pt' for padding
  \node[#2, inner xsep=6pt] (blk) at (#4,#1) {#5};
  
  % 2. Render the ghost node at the previous x (#3)
  % We use \phantom{#5} so it matches the width/height of the label exactly,
  % but remains completely invisible (empty box)
  \node[ghost, inner xsep=6pt] at (#3,#1) {\phantom{#5}};
  
  % Spring above
  \draw[hsp] (#3, #1+0.62) -- (#4, #1+0.62);
  
  % Anchor lines
  \draw[line width=0.35pt, mnnwInk] (#3, #1+0.275) -- (#3, #1+0.62);
  \draw[line width=0.35pt, mnnwInk] (#4, #1+0.275) -- (#4, #1+0.62);

  % Damper below
  \draw[line width=0.35pt, mnnwInk] (#3, #1-0.275) -- (#3, #1-0.62);
  \draw[line width=0.35pt, mnnwInk] (#4, #1-0.275) -- (#4, #1-0.62);
  
  \hdamperLong{#3}{#4}{#1-0.62}

  % DYNAMIC LABELS (K, C, and f):
  % Logic: If block moves right (#4>#3), place labels on the left (x=#3-0.2, east anchor).
  %        If block moves left (#4<#3), place labels on the right (x=#3+0.2, west anchor).
  \pgfmathsetmacro{\lblSide}{#4 > #3 ? #3 : #3}
  \pgfmathsetmacro{\lblAnch}{#4 > #3 ? "east" : "west"}

  \node[lbl, anchor=\lblAnch] at (\lblSide, #1+0.62) {#6}; % K Label
  \node[lbl, anchor=\lblAnch] at (\lblSide, #1-0.62) {#7}; % C Label
  
  % f label placed at the previous_x axis (#3) but aligned to the same side
  % y-offset of -0.2 matches your deviation arrow's height
  \pgfmathsetmacro{\fSide}{#4 > #3 ? #4 : #4}
  \pgfmathsetmacro{\fAnch}{#4 > #3 ? "west" : "east"}
  \node[ftn, anchor=\fAnch] at (\fSide, #1-0.525) {#8};

  % Deviation arrow (line remains, but label is now at the #3 axis)
  \draw[darr] (#3, #1-0.275) -- (#4, #1-0.275);

  % Update trackers for the next row
  \pgfmathsetmacro{\newSpringY}{#1+0.62}
  \xdef\lastSpringY{\newSpringY}
  \pgfmathsetmacro{\newX}{#4}
  \xdef\lastX{\newX}
}

% =====================================================================
%  GEOMETRY — full stack
% =====================================================================
\def\xLeft{-8.0}\def\xRight{8.0}
\def\xPilL{-7.4}\def\xPilR{7.4}

% Vertical levels (cm). Tight rhythm of 1.6 cm per body row.
% Wafer chain (bodies 2..7 → 6 rows), then F2, then optics (2..4 → 3 rows),
% but per user spec: optics chain attaches around F2 (lens region).
% Layout chosen:
%   F0 floor                  y = 0
%   F1 base frame             y = 1.4
%   Wafer chain  i=2..7       y = 3.2, 4.8, 6.4, 8.0, 9.6, 11.2
%   F2 metrology frame        y = 13.0   (with lens on it)
%   Optics chain  i=2..4      y = 14.4, 15.8, 17.2  (downstream of F2)
%   Reticle chain i=2..7      y = 18.8, 20.4, 22.0, 23.6, 25.2, 26.8
% =====================================================================
\def\yFloor{0}
\def\yBase{1.4}
% Wafer
\def\yWa{3.4}\def\yWb{5.4}\def\yWc{7.4}\def\yWd{9.4}\def\yWe{11.4}\def\yWf{13.4}
% Metrology
% \def\yMet{15.4}
% Optics
\def\yOa{17.2}\def\yOb{19.2}\def\yOc{21.2}
% Reticle
\def\yRa{25.2}\def\yRb{27.2}\def\yRc{29.2}\def\yRd{31.2}\def\yRe{33.2}\def\yRf{35.2}
\def\yTop{35.5}
\def\yShA{15.4}\def\yShB{23.2}

% =====================================================================
%  FLOOR  +  PILLARS  +  Pd
% =====================================================================
\filldraw[hatchfloor, draw=mnnwInk, line width=0.7pt]
  (\xRight, 0) -- (\xLeft-0.6, 0) -- (\xLeft-0.6, \yWb) -- (\xLeft-1.1, \yWb) 
  -- (\xLeft-1.1, -0.5) -- (\xRight, -0.5) -- cycle;
\node[lbl, anchor=west] at (\xRight+0.08,0.20) {$\mathcal{F}_0$};
\node[lbl, anchor=east, text=mnnwMute]
  at (\xLeft-1.1-0.08,-0.18) {Floor};

% \fill[pillar] (\xPilR,0)      rectangle (\xPilR+0.16,\yTop);

\draw[axis] (0,-0.05) -- (0,\yTop);
% \node[font=\small\itshape, anchor=south] at (0,\yTop+0.45) {$P_d$};

% =====================================================================
%  BASE FRAME  F_1
% =====================================================================
\filldraw[draw=frameEdge, line width=0.6pt, fill=frameFill]
  (7.8, \yBase-0.4) -- (-7.8, \yBase-0.4) -- (-7.8, \yTop) -- (-7.4, \yTop)
  -- (-7.4, \yShB+0.4) -- (7.8, \yShB+0.4) -- (7.8, \yShB-0.4) -- (-7.4, \yShB-0.4)
  -- (-7.4, \yShA+0.4) -- (7.8, \yShA+0.4) -- (7.8, \yShA-0.4) -- (-7.4, \yShA-0.4)
  -- (-7.4, \yBase+0.4) -- (7.8, \yBase+0.4) -- cycle;
\node[font=\small\bfseries, text=mnnwInk] at (0, \yBase) {Base Frame $\;(\mathcal{F}_1)$};

% K_1 / C_1  horizontal floor → base frame
\def\lblX{\xLeft-1.2} 

% 1. The Spring (K1)
\hdamperLong{\xLeft-0.6}{-7.8}{2.5}
\node[lbl, above=2pt] at ({(\xLeft-0.6-7.8)/2}, 2.5) {$C_1$};

% 2. The Spring (K1) - Placed on bottom
\draw[hsp] (\xLeft-0.6, 2.0) -- (-7.8, 2.0)
  node[midway, below=4pt, lbl] {$K_1$};

% =====================================================================
%  WAFER CHAIN  (j=3) — light red, blocks shifted RIGHT of Pd
%  i = 2 LS, 3 LS-flex, 4 MS, 5 MS-flex, 6 SS, 7 wafer (N3)
% =====================================================================
\pgfmathsetmacro{\shelfW}{\yBase+0.4}
\xdef\lastSpringY{\shelfW} \xdef\lastX{0}
\chainrow{\yWa}{wblk}{0}{ 0.9}{\textbf{Wafer LS} $(i_3{=}2)$}{$K_{2,3}$}{$C_{2,3}$}{$f_{2,3}$}
\chainrow{\yWb}{wflx}{0.9}{ 1.8}{flex coupling $(i_3{=}3)$}{$K_{3,3}$}{$C_{3,3}$}{$f_{3,3}$}
\chainrow{\yWc}{wblk}{1.8}{ 2.7}{\textbf{Wafer SS} $(i_3{=}4)$}{$K_{4,3}$}{$C_{4,3}$}{$f_{4,3}$}
\chainrow{\yWd}{wflx}{2.7}{ 3.6}{flex coupling $(i_3{=}5)$}{$K_{5,3}$}{$C_{5,3}$}{$f_{5,3}$}
\chainrow{\yWe}{wblk}{3.6}{ 4.5}{\textbf{Wafer Stage} $(i_3{=}6)$}{$K_{6,3}$}{$C_{6,3}$}{$f_{6,3}$}
\chainrow{\yWf}{wblk}{4.5}{ 5.4}{\textbf{Die} $(N_3{=}7)$}{$K_{N_3,3}$}{$C_{N_3,3}$}{$f_{N_3,3}$}

% =====================================================================
%  METROLOGY FRAME  F_2  +  Lens body
% =====================================================================
% \node[rfrm, minimum width=13cm, minimum height=1.0cm] (MF) at (0,\yMet)
%   {Metrology Frame $\;(\mathcal{F}_2)$};
% \node[draw=mnnwInk, line width=0.4pt, fill=lensFill,
%       minimum width=2.4cm, minimum height=0.40cm,
%       font=\scriptsize, anchor=south] at (0,\yMet+0.55) {Lens cell};
% \draw[line width=0.3pt, mnnwInk] (-0.40,\yMet+0.55) -- (-0.75,\yMet+1.05);
% \draw[line width=0.3pt, mnnwInk] ( 0.40,\yMet+0.55) -- ( 0.75,\yMet+1.05);
% \draw[line width=0.3pt, mnnwInk] (-0.40,\yMet-0.10) -- (-0.75,\yMet-0.65);
% \draw[line width=0.3pt, mnnwInk] ( 0.40,\yMet-0.10) -- ( 0.75,\yMet-0.65);

% K_2 / C_2  horizontal pillar → metrology (left side)
% \draw[hsp] (-7.4, \yMet+0.25) -- (-6.5, \yMet+0.25)
%   node[midway, above=2pt, lbl] {$K_{2,2}$};
% \hdamperLong{-7.4}{-6.5}{\yMet-0.25}
% \node[lbl, below] at (-6.95, \yMet-0.38) {$C_{2,2}$};

% =====================================================================
%  OPTICS CHAIN  (j=2)  — light purple, blocks shifted LEFT of Pd
%  Modeled as a 4-body chain whose first body (i=2) attaches to F2.
%  i = 2 lens-frame coarse, i = 3 flex, i = 4 lens cell (N2).
%  Drawn UPWARD from F2 toward the optical axis region.
% =====================================================================
\pgfmathsetmacro{\shelfO}{\yShA+0.4}
\xdef\lastSpringY{\shelfO} \xdef\lastX{0}
\chainrow{\yOa}{oblk}{0}{-0.9}{\textbf{Metrology Frame} $(i_2{=}2)$}{$K_{2,2}$}{$C_{2,2}$}{$f_{2,2}$}
\chainrow{\yOb}{oblk}{-0.9}{-1.8}{Optics Box $(i_2{=}3)$}{$K_{3,2}$}{$C_{3,2}$}{$f_{3,2}$}
\chainrow{\yOc}{oblk}{-1.8}{-2.7}{Lens Element $(i_2{=}4)$}{$K_{4,2}$}{$C_{4,2}$}{$f_{4,2}$}

% =====================================================================
%  RETICLE CHAIN  (j=1) — light blue, blocks shifted LEFT of Pd
%  i = 2 LS, 3 LS-flex, 4 MS, 5 MS-flex, 6 SS, 7 mask (N1)
% =====================================================================
\pgfmathsetmacro{\shelfR}{\yShB+0.4}
\xdef\lastSpringY{\shelfR} \xdef\lastX{0}
\chainrow{\yRa}{rblk}{0}{-0.9}{\textbf{Reticle LS} $(i_1{=}2)$}{$K_{2,1}$}{$C_{2,1}$}{$f_{2,1}$}
\chainrow{\yRb}{rflx}{-0.9}{-1.8}{flex coupling $(i_1{=}3)$}{$K_{3,1}$}{$C_{3,1}$}{$f_{3,1}$}
\chainrow{\yRc}{rblk}{-1.8}{-2.7}{\textbf{Reticle SS} $(i_1{=}4)$}{$K_{4,1}$}{$C_{4,1}$}{$f_{4,1}$}
\chainrow{\yRd}{rflx}{-2.7}{-3.6}{flex coupling $(i_1{=}5)$}{$K_{5,1}$}{$C_{5,1}$}{$f_{5,1}$}
\chainrow{\yRe}{rblk}{-3.6}{-4.5}{\textbf{Reticle Stage} $(i_1{=}6)$}{$K_{6,1}$}{$C_{6,1}$}{$f_{6,1}$}
\chainrow{\yRf}{rblk}{-4.5}{-5.4}{\textbf{Mask} $(N_1{=}7)$}{$K_{N_1,1}$}{$C_{N_1,1}$}{$f_{N_1,1}$}

% =====================================================================
%  CHAIN BRACES  (right edge — quick orienteering)
% =====================================================================
\draw[line width=0.4pt, waferEdge]   (\xPilR+0.30,\yWa-0.30) -- (\xPilR+0.30,\yWf+0.30);
\node[font=\scriptsize\bfseries, text=waferEdge, rotate=-90, anchor=south]
  at (\xPilR+0.55,{(\yWa+\yWf)/2}) {Wafer chain $(j{=}3)$};
\draw[line width=0.4pt, opticsEdge]  (\xPilR+0.30,\yOa-0.30) -- (\xPilR+0.30,\yOc+0.30);
\node[font=\scriptsize\bfseries, text=opticsEdge, rotate=-90, anchor=south]
  at (\xPilR+0.55,{(\yOa+\yOc)/2}) {Optics chain $(j{=}2)$};
\draw[line width=0.4pt, reticleEdge] (\xPilR+0.30,\yRa-0.30) -- (\xPilR+0.30,\yRf+0.30);
\node[font=\scriptsize\bfseries, text=reticleEdge, rotate=-90, anchor=south]
  at (\xPilR+0.55,{(\yRa+\yRf)/2}) {Reticle chain $(j{=}1)$};

\end{tikzpicture}
\caption{Lumped mass-spring-damper system of the three chains vertically stacked on the common base frame}
\label{fig:litho_model_full}
\end{figure}

\section{Dynamics}\label{sec:model_dynamics}

Two distinct joint types arise: \emph{actuated} stages ($i \in \{2,4\}$), which are driven by servo motors and carry a near-infinite effective stiffness that keeps $f_{i_j}$ locked to the nominal reference; and \emph{flex} stages ($i \in \{3,5,6,7\}$), which are passive compliant connections whose stiffness $K_{i_j}$ and damping $C_{i_j}$ are physical material properties of the mechanical coupling.

The full 6-DOF model is here restricted to in-plane motions\cite{Al-Rawashdeh2023_Trajectories,Dong2023_SCurve}, since rotations about the $x$ and $y$ axis, as well as translation on $z$ are used mainly for focusing. The generalized coordinate used to describe the configuration of the $i$th stage within the $j$th chain is
\begin{equation}
    q_{i_j} = \begin{bmatrix} f_{i_j} \\ \theta_{z_{i_j}} \end{bmatrix} \in \mathbb{R}^{3\times1}
\end{equation}
where $f_{i_j} = [f_{x,i_j},\, f_{y,i_j}]^T \in \mathbb{R}^{2}$ collects the in-plane translational displacements of stage $i$ relative to its parent body in chain $j$, and $\theta_{z_{i_j}} \in \mathbb{R}$ is the corresponding rotation about the out-of-plane $z$-axis. The joints between bodies within a chain are modeled using a spring-damper system~\cite{Al-Rawashdeh2022_Caracterization}. The associated generalized mass-inertia matrix for each body is
\begin{equation}
    M_{i_j} = \begin{bmatrix} m_{i_j}\,\mathbf{I}_2 & 0 \\ 0 & I_{z_{i_j}} \end{bmatrix} \in \mathbb{R}^{3\times3}
\end{equation}
and the generalized stiffness and damping matrices $K_{i_j},\, C_{i_j} \in \mathbb{R}^{3\times3}$ are block-diagonal, with a $2\times2$ translational block and a scalar rotational entry.

The absolute position $f^{0}_{i_j}$ of body $i$ is expressed as the recursive sum of its relative displacement from its parent body along the chain from the floor frame $\mathcal{F}_{0}$. Thus, for bodies with $i \geq 2$ on any chain $j \in \{1,2,3\}$:
\begin{equation} \label{eq:frameref}
    f^{0}_{i_j} = f^{0}_{(i-1)_j} + f_{i_j}
\end{equation}

Since the rotation about the $z$-axis is a scalar, the angular position recurrence is simply additive:
\begin{equation} \label{eq:angleref}
    \theta^{0}_{z_{i_j}} = \theta^{0}_{z_{(i-1)_j}} + \theta_{z_{i_j}}
\end{equation}

From this, the translational velocity expression appears as:
\begin{equation}
    \dot{f}^{0}_{i_j} = \dot{f}^{0}_{(i-1)_j} + \dot{f}_{i_j} + \omega_{(i-1)_j}
    \times
    f_{i_j}
\end{equation}

Where $\omega_{(i-1)_j}$ is the angular velocity~\cite{Ginsberg1984_AdvancedDynamics} of $\mathcal{F}_{(i-1)_j}$ with respect to the floor, $\mathcal{F}_{0}$.

Taking the second derivative gives:
\begin{align}
    \begin{split}
        \ddot{f}^{0}_{i_j} &= \ddot{f}^{0}_{(i-1)_j} + \ddot{f}_{i_j}
        + \omega_{(i-1)_j} \times \dot{f}_{i_j}
        + \dot{\omega_{(i-1)_j}} \times f_{i_j}
        + \omega_{(i-1)_j} \times \left[
            \dot{f}_{i_j} + \omega_{(i-1)_j} \times f_{i_j}
        \right] \\
        &= \ddot{f}^{0}_{(i-1)_j} + \ddot{f}_{i_j}
        + \omega_{(i-1)_j} \times \dot{f}_{i_j}
        + \dot{\omega_{(i-1)_j}} \times f_{i_j}
        + \omega_{(i-1)_j} \times \dot{f}_{i_j}
        + \omega_{(i-1)_j} \times \left[
            \omega_{(i-1)_j} \times f_{i_j}
        \right] \\
        &= \ddot{f}^{0}_{(i-1)_j} + \ddot{f}_{i_j}
        + 2 \left[ \omega_{(i-1)_j} \times \dot{f}_{i_j} \right]
        + \dot{\omega_{(i-1)_j}} \times f_{i_j}
        + \omega_{(i-1)_j} \times \left[
            \omega_{(i-1)_j} \times f_{i_j}
        \right]
    \end{split}
\end{align}

Since we are concerned with planar movements in this work, we restrict the angular velocity of each body purely about the $z$-axis as $\omega = \dot{\theta}_z\,\hat{z}$, so that:
\begin{equation}
    \omega_{(i-1)_j} = \dot{\theta}^{0}_{z_{(i-1)_j}}\,\hat{z} \notag
\end{equation}

We substitute the cross-products as matrix-vector products using the skew-symmetric matrix. In the 6-DOF model~\cite{Al-Rawashdeh2022_Caracterization}, it appears as:
\begin{equation}
    [\omega]
    =
    \begin{bmatrix}
        0 & -\omega_z & \omega_y \\
        \omega_z & 0 & -\omega_x \\
        -\omega_y & \omega_x & 0
    \end{bmatrix}
    \notag
\end{equation}
Clearly, here we can set $\omega_x = 0$ and $\omega_y = 0$, thus $\omega = [0,\, 0,\, \dot{\theta}_z]^T$ and the matrix becomes:
\begin{equation}
    [\omega]
    =
    \begin{bmatrix}
        0 & -\omega_z & 0 \\
        \omega_z & 0 & 0 \\
        0 & 0 & 0
    \end{bmatrix}
    \notag
\end{equation}

Considering our 2D displacement vector, we only take the upper-left $2\times2$ 
sub-matrix, and factoring out the scalar $\dot{\theta}^{0}_{z_{(i-1)_j}}$:
\begin{equation} \label{eq:skew}
    [\omega_{(i-1)_j}] = \dot{\theta}^{0}_{z_{(i-1)_j}}
    \begin{bmatrix} 0 & -1 \\ 1 & 0 \end{bmatrix}
\end{equation}
such that $\omega_{(i-1)_j} \times f_{i_j} = [\omega_{(i-1)_j}]\,f_{i_j}$. 
Applying this matrix twice gives the centripetal term:
\begin{equation}
    [\omega_{(i-1)_j}]^2 = -\left(\dot{\theta}^{0}_{z_{(i-1)_j}}\right)^2 \mathbf{I}_2
\end{equation}
so that $\omega_{(i-1)_j} \times [\omega_{(i-1)_j} \times f_{i_j}] =
[\omega_{(i-1)_j}]^2\,f_{i_j}$. Substituting and expanding
\eqref{eq:frameref} recursively from the floor frame, the translational kinematics become:
\begin{align}
    f^{0}_{i_j}
        &= f^{0}_{1} + \sum^{i}_{k=2} f_{k_j}
        \label{eq:pos_expanded} \\
    \dot{f}^{0}_{i_j}
        &= \dot{f}^{0}_{1} + \sum^{i}_{k=2} \dot{f}_{k_j}
         + \sum^{i}_{k=2} [\omega_{(k-1)_j}]\, f_{k_j}
        \label{eq:vel_expanded} \\
    \ddot{f}^{0}_{i_j}
        &= \ddot{f}^{0}_{1} + \sum^{i}_{k=2} \ddot{f}_{k_j}
         + \sum^{i}_{k=2} \left[
             2\,[\omega_{(k-1)_j}]\,\dot{f}_{k_j}
             + [\dot{\omega}_{(k-1)_j}]\,f_{k_j}
             + [\omega_{(k-1)_j}]^2\,f_{k_j}
         \right]
        \label{eq:acc_expanded}
\end{align}

Since the rotation $\theta_{z_{i_j}}$ is a scalar, no Coriolis or centripetal cross-terms arise in the angular kinematics. Expanding~\eqref{eq:angleref} recursively from the floor frame gives the rotational analogues:
\begin{align}
    \theta^{0}_{z_{i_j}}
        &= \theta^{0}_{z_1} + \sum^{i}_{k=2} \theta_{z_{k_j}}
        \label{eq:angle_expanded} \\
    \dot{\theta}^{0}_{z_{i_j}}
        &= \dot{\theta}^{0}_{z_1} + \sum^{i}_{k=2} \dot{\theta}_{z_{k_j}}
        \label{eq:angvel_expanded} \\
    \ddot{\theta}^{0}_{z_{i_j}}
        &= \ddot{\theta}^{0}_{z_1} + \sum^{i}_{k=2} \ddot{\theta}_{z_{k_j}}
        \label{eq:angacc_expanded}
\end{align}

Thus the equation of the angular motion of the $i$th body in the $j$th chain w.r.t $\mathcal{F}_{0}$ is:
\begin{equation}
    {}^{0}\!I_j \, \dot{\omega}_j + \left( [\omega_j] \right) {}^{0}\!I_j \, \omega_j  = m_j^0
\end{equation}

Under the planar restriction $\omega_x = \omega_y = 0$ the gyroscopic term vanishes like:
\begin{equation}
    \left([\omega_{i_j}]\,{}^0\!I_{i_j}\right)\omega_{i_j}
    =
    \begin{bmatrix}
        0 & -\dot{\theta}_z & 0 \\
        \dot{\theta}_z & 0 & 0 \\
        0 & 0 & 0
    \end{bmatrix}
    \begin{bmatrix}
        I_{xx} & 0 & 0 \\
        0 & I_{yy} & 0 \\
        0 & 0 & I_{zz}
    \end{bmatrix}
    \begin{bmatrix}
        0 \\
        0 \\
        \dot{\theta}_z
    \end{bmatrix}
    = 0
    \notag
\end{equation}

Leaving only:
\begin{equation}
    I_{z_{i_j}}\ddot{\theta}^{0}_{z_{i_j}} = m^0_{z_{i_j}}
\end{equation}

where $I_{z_{i_j}}$ is the aggregated moment of inertia of all bodies from $i$ to the end-effector $N_j$, referred to the $z$-axis of body $i$ via Steiner's theorem, and $m^0_{z_{i_j}} \in \mathbb{R}$ is the net applied moment acting on body $i$ in chain $j$ about the $z$-axis.

\paragraph{Steiner's theorem and inertia aggregation.}
The parallel axis (Steiner's) theorem states that the moment of inertia of a rigid body about any axis parallel to one through its center of mass (CoG) satisfies
\begin{equation} \label{eq:steiner}
    I_z = \tilde{I}_z + m\,\|f\|^2,
\end{equation}
where $\tilde{I}_z$ is the body's own moment of inertia about the $z$-axis through its CoG, $m$ is its mass, and $\|f\|^2 = f_x^2 + f_y^2$ is the squared in-plane distance between the two parallel axes. In a serial chain, the inertia of every downstream body must be shifted to the frame of the current body using~\eqref{eq:steiner}. Following the characterization article\cite{Al-Rawashdeh2022_Caracterization} (Algorithm~1 therein), the aggregated inertia for the $j$th chain is computed recursively from the end-effector inward. In the planar restriction, the full $3\times3$ inertia matrix of the original algorithm reduces to a scalar $z$-axis inertia:

\begin{algorithm}[H]
\caption{Planar inertia aggregation for the $j$th chain — planar reduction of Algorithm~1 in~\cite{Al-Rawashdeh2022_Caracterization}}
\begin{algorithmic}[1]
\Require{CoG inertias $\tilde{I}_{z_{i_j}}$, masses $m_{i_j}$, relative positions $f_{k_j}$ for $i = 2,\ldots,N_j$}
\Ensure{Aggregated inertias $^0I_{z_{i_j}}$ for $i = 2,\ldots,N_j$}
\State \textbf{Initialize:} $^0I_{z_{N_j}} \leftarrow \tilde{I}_{z_{N_j}}$
\If{$2 \leq i < N_j$}
    \For{$k = N_j,\;\ldots,\;i+1$}
        \State $\tilde{m}_{(k-1)_j} \leftarrow \displaystyle\sum_{b=k}^{N_j} m_{b_j}$
        \Comment{accumulated downstream mass}
        \State $^0I_{z_{(k-1)_j}} \leftarrow {}^0I_{z_{k_j}} + \tilde{I}_{z_{(k-1)_j}} + \tilde{m}_{(k-1)_j}\bigl(f_{x,k_j}^2 + f_{y,k_j}^2\bigr)$
        \Comment{Steiner shift from frame $k$ to frame $k\!-\!1$}
    \EndFor
\EndIf
\end{algorithmic}
\end{algorithm}

Here $\tilde{m}_{(k-1)_j} = \sum_{b=k}^{N_j} m_{b_j}$ is the total mass of all downstream bodies (from $k$ to the end-effector), $\tilde{I}_{z_{(k-1)_j}}$ is the local inertia of body $k-1$ about its own CoG, and $\tilde{m}_{(k-1)_j}(f_{x,k_j}^2 + f_{y,k_j}^2)$ is the Steiner correction that shifts the accumulated inertia from body $k$'s CoG to body $(k-1)$'s CoG. By the time the algorithm terminates, $^0I_{z_{2_j}}$ contains the inertia of the entire chain $j$ referred to the CoG of its first body.

The aggregated inertia of the base frame (body~1) applies the same shift one more time, from body~2's CoG to the base frame:
\begin{equation} \label{eq:Iz0}
    {}^0I_{z_1} = \tilde{I}_{z_1}
    + \sum_{j=1}^{3}\!\Bigl({}^0I_{z_{2_j}}
      + \tilde{m}_{1_j}\bigl(f_{x,2_j}^2 + f_{y,2_j}^2\bigr)\Bigr),
\end{equation}
where $\tilde{m}_{1_j} = \sum_{b=2}^{N_j} m_{b_j}$ is the total mass of chain $j$. This is the quantity denoted $I_{z_0}$ in~\eqref{eq:baseframerot}.

Expanding $\ddot{\theta}^0_{z_{i_j}}$ using~\eqref{eq:angacc_expanded}:
\begin{equation} \label{eq:torqueeq_expanded}
    I_{z_{i_j}}\ddot{\theta}^{0}_{z_1} 
    + I_{z_{i_j}}\sum^{i}_{k=2}\ddot{\theta}_{z_{k_j}}
    = m^0_{z_{i_j}}
\end{equation}

\subsection{For the base frame}

To write the equations that govern the base frame, we start from Newton's equation $F = ma$. For every body ($i \geq 2$) within chain $j$, it connects to body $i-1$ through a spring above and body $i+1$ through the spring below:
\begin{align} \label{eq:forceeq}
    \begin{split}
        m_{i_j}\ddot{f}^{0}_{1}
        + m_{i_j}\sum^{i}_{k=2}\ddot{f}_{k_j}
        + m_{i_j}\sum^{i}_{k=2}
        &\left[
            2\,[\omega_{(k-1)_j}]\,\dot{f}_{k_j}
            + [\dot{\omega}_{(k-1)_j}]\,f_{k_j}
            + [\omega_{(k-1)_j}]^2\,f_{k_j}
        \right] \\
        - K_i\left( f^0_{{(i-1)}_j} - f^0_{{i}_j}\right)
        - C_i\left( \dot{f}^0_{{(i-1)}_j} - \dot{f}^0_{{i}_j}\right)
        &+ K_{i+1}\left( f^0_{{i}_j} - f^0_{{(i+1)}_j}\right)
        + C_{i+1}\left( \dot{f}^0_{{i}_j} - \dot{f}^0_{{(i+1)}_j}\right)
        = u^0_{i_j}
    \end{split}
\end{align}

Summing~\eqref{eq:forceeq} over all bodies in a chain (taking the optics chain $j=2$, bodies $i=2,3,4$ as a representative example), the internal spring forces cancel by Newton's third law:
\begin{align}
    &- K_2(f^0_1 - f^0_2) - C_2(\dot{f}^0_1 - \dot{f}^0_2) \notag \\
    &\quad \cancelto{0}{
        + K_3(f^0_2 - f^0_3) + C_3(\dot{f}^0_2 - \dot{f}^0_3)
        - K_3(f^0_2 - f^0_3) - C_3(\dot{f}^0_2 - \dot{f}^0_3)
    } \notag \\
    &\quad \cancelto{0}{
        + K_4(f^0_3 - f^0_4) + C_4(\dot{f}^0_3 - \dot{f}^0_4)
        - K_4(f^0_3 - f^0_4) - C_4(\dot{f}^0_3 - \dot{f}^0_4)
    } \notag
\end{align}

Only the boundary term $- K_2(f^0_1 - f^0_2) - C_2(\dot{f}^0_1 - \dot{f}^0_2)$ survives, which corresponds to the connection between body $i = 2$ and the base frame. This extends to any $N_j$-body chain. Recalling~\eqref{eq:frameref}, $f^{0}_{1} - f^{0}_{2_j} = -f_{2_j}$, and following~\cite{Al-Rawashdeh2022_Caracterization} with ${}^0K_{i_j}={}^0R_{i_j} K_{i_j} {}^0R^T_{i_j}$ and ${}^0C_{i_j}={}^0R_{i_j} C_{i_j} {}^0R^T_{i_j}$, the total translational equation of the base frame (absorbing the balance masses and collecting all chain contributions) is:
\begin{align} \label{eq:baseframeeq}
    \begin{split}
        m_0\,\ddot{f}^{0}_{1}
        + {}^0\bar{C}_1\,\dot{f}_1
        + {}^0\bar{K}_1\,f_1
        - \sum_{j=1}^{3}\left\{
            {}^0\hat{C}_{2_j}\,\dot{f}_{2_j} + {}^0\hat{K}_{2_j}\,f_{2_j}
        \right\}
        &= u^0_{1}
    \end{split}
\end{align}
where $m_0 = m_{1} + \sum_{j=1}^{3}\!\left(m^j_{B} + \sum_{i=2}^{N_j} m_{i_j}\right)$ is the total system mass, and the stiffness and damping matrices denote the $2\times2$ translational block. Following~\cite{Al-Rawashdeh2022_Caracterization}, ${}^0\bar{K}_1 = {}^0K_1$ and ${}^0\bar{C}_1 = {}^0C_1$, while the hatted chain couplings are ${}^0\hat{C}_{2_j} = {}^0C_{2_j}$ and ${}^0\hat{K}_{2_j} = {}^0K_{2_j} + {}^0C_{2_j}[\omega_1]$; the accents are defined in general in~\eqref{eq:accented} below.

By the exact same cancellation argument applied to the moment balance about the $z$-axis — the moment exerted by the spring connecting body $i$ to body $i-1$ is $-K_{i_j}\,\theta_{z_{i_j}}$, so internal moments cancel identically — the rotational equation of the base frame is:
\begin{align} \label{eq:baseframerot}
    \begin{split}
        I_{z_0}\,\ddot{\theta}^{0}_{z_1}
        + {}^0\bar{C}_1\,\dot{\theta}_{z_1}
        + {}^0\bar{K}_1\,\theta_{z_1}
        - \sum_{j=1}^{3}\left\{
            {}^0\hat{C}_{2_j}\,\dot{\theta}_{z_{2_j}}
            + {}^0\hat{K}_{2_j}\,\theta_{z_{2_j}}
        \right\}
        &= \tau^0_{1}
    \end{split}
\end{align}
where $I_{z_0} = I_{z_1} + \sum_{j=1}^{3}\!\left(I^j_{z_B} + \sum_{i=2}^{N_j} I_{z_{i_j}}\right)$ is the total moment of inertia of the system about the $z$-axis, ${}^0\bar{C}_{i_j}$ and ${}^0\bar{K}_{i_j}$ (and their hatted counterparts, which coincide with them for the scalar rotational entries) are the scalar rotational damping and stiffness of the joint connecting body $i$ to its predecessor in chain $j$, and $\tau^0_1$ is the net applied torque on the base frame.

\subsection{For each chain}

Recalling~\eqref{eq:forceeq} and~\eqref{eq:torqueeq_expanded}, for the first body in the chain ($i=2$), which connects to the base frame above and to body $i=3$ below, the translational equation is:
\begin{equation}
        \bar{m}_{2_j}
        \ddot{f}^{0}_{2_j}
        + {}^0\bar{C}_{2_j}\,\dot{f}_{2_j}
        + {}^0\bar{K}_{2_j}\,f_{2_j}
        - \left\{
            {}^0\hat{C}_{3_j}\,\dot{f}_{3_j} + {}^0\hat{K}_{3_j}\,f_{3_j}
        \right\}
        = u^0_{2_j}
\end{equation}
and the rotational equation is:
\begin{equation}
        I_{z_{2_j}}\ddot{\theta}^{0}_{z_{2_j}}
        + {}^0\bar{C}_{2_j}\,\dot{\theta}_{z_{2_j}}
        + {}^0\bar{K}_{2_j}\,\theta_{z_{2_j}}
        - {}^0\hat{C}_{3_j}\,\dot{\theta}_{z_{3_j}}
        - {}^0\hat{K}_{3_j}\,\theta_{z_{3_j}}
        = \tau^0_{2_j}
\end{equation}

For bodies $3\leq i< N_j$, the same structure holds with bodies $2, \dots, i-1$ in between. The translational equation is:
\begin{equation}
        \bar{m}_{i_j}
        \ddot{f}^{0}_{i_j}
        + {}^0\bar{C}_{i_j}\,\dot{f}_{i_j}
        + {}^0\bar{K}_{i_j}\,f_{i_j}
        -
        \left\{
            {}^0\hat{C}_{(i+1)_j}\,\dot{f}_{(i+1)_j} + {}^0\hat{K}_{(i+1)_j}\,f_{(i+1)_j}
        \right\}
        = u^0_{i_j}
\end{equation}
and the rotational equation is:
\begin{equation}
        I_{z_{i_j}}\ddot{\theta}^{0}_{z_{i_j}}
        + {}^0\bar{C}_{i_j}\,\dot{\theta}_{z_{i_j}}
        + {}^0\bar{K}_{i_j}\,\theta_{z_{i_j}}
        - {}^0\hat{C}_{(i+1)_j}\,\dot{\theta}_{z_{(i+1)_j}}
        - {}^0\hat{K}_{(i+1)_j}\,\theta_{z_{(i+1)_j}}
        = \tau^0_{i_j}
\end{equation}

Finally, for body $i = N_j$ (the end effector), there is no subsequent body in the chain. The translational and rotational equations are:
\begin{equation}
        \bar{m}_{{N_j}_j}
        \ddot{f}^{0}_{{N_j}_j}
        + {}^0\bar{C}_{{N_j}_j}\,\dot{f}_{{N_j}_j}
        + {}^0\bar{K}_{{N_j}_j}\,f_{{N_j}_j}
        = 0
\end{equation}
\begin{equation}
        I_{z_{{N_j}_j}}\ddot{\theta}^{0}_{z_{{N_j}_j}}
        + {}^0\bar{C}_{{N_j}_j}\,\dot{\theta}_{z_{{N_j}_j}}
        + {}^0\bar{K}_{{N_j}_j}\,\theta_{z_{{N_j}_j}}
        = 0
\end{equation}

\subsection{Accented matrices and the nominal condition}

Following Appendix~A of~\cite{Al-Rawashdeh2022_Caracterization}, the accented matrices appearing in the equations above are not merely rotated joint matrices: they absorb the rotational cross-terms of~\eqref{eq:acc_expanded} into effective damping and stiffness. With $[\omega]$ the planar skew form of~\eqref{eq:skew} and $B_{i_j} = [\dot{\omega}_{i_j}] + [\omega_{i_j}]^2$:
\begin{align}
    \bar{m}_{i_j} &= \sum_{k=i}^{N_j} m_{k_j} \notag \\
    {}^0\bar{C}_{i_j} &= {}^0C_{i_j} + 2m_{i_j}\,[\omega_{(i-1)_j}] \notag \\
    {}^0\bar{K}_{i_j} &= {}^0K_{i_j} + {}^0C_{i_j}\,[\omega_{(i-1)_j}] + m_{i_j}\,B_{(i-1)_j} \label{eq:accented} \\
    {}^0\hat{C}_{(i+1)_j} &= {}^0C_{(i+1)_j} \notag \\
    {}^0\hat{K}_{(i+1)_j} &= {}^0K_{(i+1)_j} + {}^0C_{(i+1)_j}\,[\omega_{i_j}] \notag
\end{align}
The bar thus marks the matrices acting on the body's \emph{own} joint coordinates, which collect the Coriolis ($2m[\omega]$), damper-convection ($C[\omega]$), and Euler--centripetal ($mB$) contributions, while the hat marks the coupling to the \emph{downstream} joint coordinates, which inherits only the convective damper term: the damper reacts to the inertial velocity $\dot{f} + [\omega]f$ of~\eqref{eq:vel_expanded}, so part of its force is proportional to position and migrates into the effective stiffness. The scalar rotational entries admit no such cross-terms in the planar restriction, so about the $z$-axis the accents are inert and $\hat{(\cdot)} = \bar{(\cdot)} = (\cdot)$. For $3 \leq i < N_j$ the model of~\cite{Al-Rawashdeh2022_Caracterization} additionally couples body $i$ to every upstream joint through the two-index matrices ${}^0\hat{C}_{i_j,k} = 2m_{i_j}\,[\omega_{(k-1)_j}]$ and ${}^0\hat{K}_{i_j,k} = m_{i_j}\,B_{(k-1)_j}$ for $2 \leq k < i$, omitted above.

Every one of these corrections is proportional to $\dot{\theta}_z$ or $\ddot{\theta}_z$ of a parent body. At the nominal operating condition $\dot{\theta}_z = \ddot{\theta}_z = 0$ --- the regulated step-and-scan regime studied throughout this work --- they vanish identically: the accented matrices reduce to the frame-rotated ${}^0K_{i_j}$ and ${}^0C_{i_j}$, the inertial accelerations of~\eqref{eq:acc_expanded} reduce to sums of relative accelerations, and the equations of motion coincide with Eqs.~(2.6)--(2.8) of~\cite{Al-Rawashdeh2022_Caracterization} and become linear time-invariant. Away from the nominal condition, the corrections re-enter as the parameter-varying part $A_1(\Omega,\dot{\Omega})$ of the decomposition given in~\cite{Al-Rawashdeh2022_Caracterization}. Note that the scheduling variables $\Omega$ are themselves state variables --- the $\dot{\theta}_z$ of each body --- so the cross-terms are quadratic in the state (and the rotated matrices ${}^0K_{i_j}$, ${}^0C_{i_j}$ carry a further state dependence through ${}^0R_{i_j}(\theta_z)$): the machine model is, strictly, nonlinear, a \emph{quasi}-LPV system whose frozen nominal dynamics are the LTI model above. The MuJoCo test bed used throughout this work assembles the full nonlinear multibody dynamics from this lumped topology directly, retaining these couplings without approximation; the LTI form above is recovered only as its frozen nominal dynamics at $\dot{\theta}_z = \ddot{\theta}_z = 0$, not as a separately integrated reduced model.

\subsection{State-space formulation}

Because each body possesses three planar degrees of freedom ($f_x$, $f_y$, $\theta_z$) along with their corresponding time derivatives, it contributes exactly six states to the system. Following the machine description of~\cite{Al-Rawashdeh2022_Caracterization} (Section~2.2 therein), the complete machine model comprises 6 bodies in the reticle chain ($N_1 = 7$), 3 bodies in the optics chain ($N_2 = 4$), and 6 bodies in the wafer chain ($N_3 = 7$). Including the common base frame (body 1), the system features a total of $1 + (N_1-1) + (N_2-1) + (N_3-1) = 16$ bodies. Consequently, the global state vector $x$ encompasses $6 \times 16 = 96$ states:

\begin{equation}
    x =
    \begin{bmatrix}
        \underbrace{x_1, \ldots, x_6}_{\text{base frame}} \;,\;
        \underbrace{x_7, \ldots, x_{6N_1}}_{\text{chain 1: bodies } 2 \to N_1} \;,\;
        \underbrace{\ldots}_{\text{chain 2}} \;,\;
        \underbrace{\ldots}_{\text{chain 3}}
    \end{bmatrix}^T
\end{equation}

The explicit mapping for each state is detailed in the table below. Here, $p$ denotes the global sequential index of the $i$-th body in the $j$-th chain within the concatenated list of all bodies (starting with the base frame, followed sequentially by chains 1, 2, and 3). It is crucial to note that every positional entry within the state vector represents a \emph{relative} joint displacement; the corresponding absolute inertial position $f^0_{i_j}$ can be reconstructed using the recursive summations from~\eqref{eq:pos_expanded} to~\eqref{eq:angle_expanded}.

\begin{table}[htbp]
\centering
\caption{State-space mapping for the global concatenated system vector}
\label{tab:state_mapping}
\begin{tabular}{cll}
\toprule
State-space Variable & Expression & Meaning \\
\midrule
$x_1$       & $f_{x_1}$                  & base frame $x$-position \\
$x_2$       & $f_{y_1}$                  & base frame $y$-position \\
$x_3$       & $\theta_{z_1}$             & base frame rotation \\
$x_4$       & $\dot{f}_{x_1}$            & base frame $x$-velocity \\
$x_5$       & $\dot{f}_{y_1}$            & base frame $y$-velocity \\
$x_6$       & $\dot{\theta}_{z_1}$       & base frame angular velocity \\
$x_7$       & $f_{x_{2_1}}$              & chain 1, body 2, $x$-position \\
$x_8$       & $f_{y_{2_1}}$              & chain 1, body 2, $y$-position \\
$x_9$       & $\theta_{z_{2_1}}$         & chain 1, body 2, rotation \\
$x_{10}$    & $\dot{f}_{x_{2_1}}$        & chain 1, body 2, $x$-velocity \\
$x_{11}$    & $\dot{f}_{y_{2_1}}$        & chain 1, body 2, $y$-velocity \\
$x_{12}$    & $\dot{\theta}_{z_{2_1}}$   & chain 1, body 2, angular velocity \\
$\vdots$    & $\vdots$                   & $\vdots$ \\
$x_{6p-5}$  & $f_{x_{i_j}}$              & chain $j$, body $i$, $x$-position \\
$x_{6p-4}$  & $f_{y_{i_j}}$              & chain $j$, body $i$, $y$-position \\
$x_{6p-3}$  & $\theta_{z_{i_j}}$         & chain $j$, body $i$, rotation \\
$x_{6p-2}$  & $\dot{f}_{x_{i_j}}$        & chain $j$, body $i$, $x$-velocity \\
$x_{6p-1}$  & $\dot{f}_{y_{i_j}}$        & chain $j$, body $i$, $y$-velocity \\
$x_{6p}$    & $\dot{\theta}_{z_{i_j}}$   & chain $j$, body $i$, angular velocity \\
$\vdots$    & $\vdots$                   & $\vdots$ \\
\bottomrule
\end{tabular}
\end{table}

Solving~\eqref{eq:baseframeeq} and~\eqref{eq:baseframerot} for the base-frame accelerations yields:
\begin{align}
    \ddot{f}^{0}_{1}
    &= -\frac{{}^0\bar{C}_1}{m_0}\,\dot{f}_1
       -\frac{{}^0\bar{K}_1}{m_0}\,f_1
       +\frac{1}{m_0}\sum_{j=1}^{3}\left\{
           {}^0\bar{C}_{2_j}\,\dot{f}_{2_j} + {}^0\bar{K}_{2_j}\,f_{2_j}
       \right\}
       +\frac{u^0_{1}}{m_0} \label{eq:f1ddot} \\[4pt]
    \ddot{\theta}^{0}_{z_1}
    &= -\frac{{}^0\bar{C}_1}{I_{z_0}}\,\dot{\theta}_{z_1}
       -\frac{{}^0\bar{K}_1}{I_{z_0}}\,\theta_{z_1}
       +\frac{1}{I_{z_0}}\sum_{j=1}^{3}\left\{
           {}^0\bar{C}_{2_j}\,\dot{\theta}_{z_{2_j}} + {}^0\bar{K}_{2_j}\,\theta_{z_{2_j}}
       \right\}
       +\frac{\tau^0_{1}}{I_{z_0}} \label{eq:th1ddot}
\end{align}

Expressed in the standard matrix form $\dot{x} = Ax + Bu$, the first six rows governing the base-frame dynamics are:
\setcounter{MaxMatrixCols}{20}
\begin{equation}
\begin{bmatrix}\dot{x}_1\\\dot{x}_2\\\dot{x}_3\\\dot{x}_4\\\dot{x}_5\\\dot{x}_6\end{bmatrix}
=
\begin{bmatrix}
    0 & 0 & 0 & 1 & 0 & 0 & 0 & 0 & 0 & 0 & \cdots \\
    0 & 0 & 0 & 0 & 1 & 0 & 0 & 0 & 0 & 0 & \cdots \\
    0 & 0 & 0 & 0 & 0 & 1 & 0 & 0 & 0 & 0 & \cdots \\
    -\frac{K_{1,xx}}{m_0} & -\frac{K_{1,xy}}{m_0} & 0 &
    -\frac{C_{1,xx}}{m_0} & -\frac{C_{1,xy}}{m_0} & 0 &
    \frac{K_{2_1,xx}}{m_0} & \frac{K_{2_1,xy}}{m_0} & 0 &
    \frac{C_{2_1,xx}}{m_0} & \cdots \\
    -\frac{K_{1,yx}}{m_0} & -\frac{K_{1,yy}}{m_0} & 0 &
    -\frac{C_{1,yx}}{m_0} & -\frac{C_{1,yy}}{m_0} & 0 &
    \frac{K_{2_1,yx}}{m_0} & \frac{K_{2_1,yy}}{m_0} & 0 &
    \frac{C_{2_1,yx}}{m_0} & \cdots \\
    0 & 0 & -\frac{K_{1,zz}}{I_{z_0}} & 0 & 0 &
    -\frac{C_{1,zz}}{I_{z_0}} & 0 & 0 &
    \frac{K_{2_1,zz}}{I_{z_0}} & 0 & \cdots \\
\end{bmatrix}
x
+
\begin{bmatrix}
    0 & 0 & 0 & \cdots \\
    0 & 0 & 0 & \cdots \\
    0 & 0 & 0 & \cdots \\
    \frac{1}{m_0} & 0 & 0 & \cdots \\
    0 & \frac{1}{m_0} & 0 & \cdots \\
    0 & 0 & \frac{1}{I_{z_0}} & \cdots
\end{bmatrix}
\begin{bmatrix}u^0_{x,1}\\u^0_{y,1}\\\tau^0_1\\\vdots\end{bmatrix}
\end{equation}

\subsubsection{General State-Space Form for the Chain Bodies}

To systematically construct the global state-space matrices, we must reframe the equations of motion for the chain bodies in terms of the state vector variables. We map the physical planar coordinates of the $p$-th sequential body to the local state vectors: translational position $\mathbf{x}_{2p-1} = [x_{6p-5},\, x_{6p-4}]^T$, translational velocity $\dot{\mathbf{x}}_{2p-1} = [x_{6p-2},\, x_{6p-1}]^T$, rotational position $\mathbf{x}_{2p} = x_{6p-3}$, and rotational velocity $\dot{\mathbf{x}}_{2p} = x_{6p}$.

The state derivatives are dictated by the relative accelerations between adjacent bodies. By applying the kinematic identity $\ddot{f}_{i_j} = \ddot{f}^0_{i_j} - \ddot{f}^0_{(i-1)_j}$ from~\eqref{eq:acc_expanded}, we determine the interactive forces acting on each body within its parent's reference frame. Let $p$ represent the global index of body $i$ in chain $j$, while $p-1$ and $p+1$ denote the indices of its immediate upstream and downstream neighbors respectively (with the base frame positioned at $p=1$). 


For the \textbf{first body in each chain} ($i=2$, index $p$, upstream index $1$), the translational state derivatives ($\ddot{\mathbf{x}}_{2p-1} = \ddot{f}_{2_j}$) and rotational state derivatives ($\ddot{\mathbf{x}}_{2p} = \ddot{\theta}_{z_{2_j}}$) take the form:
\begin{align}
    \ddot{\mathbf{x}}_{2p-1} &= 
    \frac{1}{\bar{m}_{2_j}} \Big[ u^0_{2_j} - {}^0\bar{C}_{2_j}\dot{\mathbf{x}}_{2p-1} - {}^0\bar{K}_{2_j}\mathbf{x}_{2p-1} + {}^0\hat{C}_{3_j}\dot{\mathbf{x}}_{2p+1} + {}^0\hat{K}_{3_j}\mathbf{x}_{2p+1} \Big] \notag \\
    &\quad - \frac{1}{m_0} \Big[ u^0_1 - {}^0\bar{C}_1\dot{\mathbf{x}}_{1} - {}^0\bar{K}_1 \mathbf{x}_{1} + \sum_{k=1}^3 \left( {}^0\hat{C}_{2_k}\dot{\mathbf{x}}_{2p_{2_k}-1} + {}^0\hat{K}_{2_k}\mathbf{x}_{2p_{2_k}-1} \right) \Big] \\
    \ddot{\mathbf{x}}_{2p} &= 
    \frac{1}{I_{z_{2_j}}} \Big[ \tau^0_{2_j} - {}^0\bar{c}_{\theta,2_j}\dot{\mathbf{x}}_{2p} - {}^0\bar{k}_{\theta,2_j}\mathbf{x}_{2p} + {}^0\hat{c}_{\theta,3_j}\dot{\mathbf{x}}_{2p+2} + {}^0\hat{k}_{\theta,3_j}\mathbf{x}_{2p+2} \Big] \notag \\
    &\quad - \frac{1}{I_{z_0}} \Big[ \tau^0_1 - {}^0\bar{C}_1\dot{\mathbf{x}}_{2} - {}^0\bar{K}_1 \mathbf{x}_{2} + \sum_{k=1}^3 \left( {}^0\hat{c}_{\theta,2_k}\dot{\mathbf{x}}_{2p_{2_k}} + {}^0\hat{k}_{\theta,2_k}\mathbf{x}_{2p_{2_k}} \right) \Big]
\end{align}

For a \textbf{general intermediate body} ($2 < i < N_j$, index $p$), the dynamics are shaped by upstream, local, and downstream interactions:
\begin{align}
    \ddot{\mathbf{x}}_{2p-1} &= 
    \frac{1}{\bar{m}_{i_j}} \Big[ u^0_{i_j} - {}^0\bar{C}_{i_j}\dot{\mathbf{x}}_{2p-1} - {}^0\bar{K}_{i_j}\mathbf{x}_{2p-1} + {}^0\hat{C}_{(i+1)_j}\dot{\mathbf{x}}_{2p+1} + {}^0\hat{K}_{(i+1)_j}\mathbf{x}_{2p+1} \Big] \notag \\
    &\quad - \frac{1}{\bar{m}_{(i-1)_j}} \Big[ u^0_{(i-1)_j} - {}^0\bar{C}_{(i-1)_j}\dot{\mathbf{x}}_{2p-3} - {}^0\bar{K}_{(i-1)_j}\mathbf{x}_{2p-3} + {}^0\hat{C}_{i_j}\dot{\mathbf{x}}_{2p-1} + {}^0\hat{K}_{i_j}\mathbf{x}_{2p-1} \Big] \\
    \ddot{\mathbf{x}}_{2p} &= 
    \frac{1}{I_{z_{i_j}}} \Big[ \tau^0_{i_j} - {}^0\bar{c}_{\theta,i_j}\dot{\mathbf{x}}_{2p} - {}^0\bar{k}_{\theta,i_j}\mathbf{x}_{2p} + {}^0\hat{c}_{\theta,(i+1)_j}\dot{\mathbf{x}}_{2p+2} + {}^0\hat{k}_{\theta,(i+1)_j}\mathbf{x}_{2p+2} \Big] \notag \\
    &\quad - \frac{1}{I_{z_{(i-1)_j}}} \Big[ \tau^0_{(i-1)_j} - {}^0\bar{c}_{\theta,(i-1)_j}\dot{\mathbf{x}}_{2p-2} - {}^0\bar{k}_{\theta,(i-1)_j}\mathbf{x}_{2p-2} + {}^0\hat{c}_{\theta,i_j}\dot{\mathbf{x}}_{2p} + {}^0\hat{k}_{\theta,i_j}\mathbf{x}_{2p} \Big]
\end{align}

Finally, for the \textbf{end-effector body} ($i = N_j$, index $p$), the downstream force terms naturally vanish as there are no subsequent bodies:
\begin{align}
    \ddot{\mathbf{x}}_{2p-1} &= 
    \frac{1}{\bar{m}_{{N_j}_j}} \Big[ u^0_{{N_j}_j} - {}^0\bar{C}_{{N_j}_j}\dot{\mathbf{x}}_{2p-1} - {}^0\bar{K}_{{N_j}_j}\mathbf{x}_{2p-1} \Big] \notag \\
    &\quad - \frac{1}{\bar{m}_{({N_j}-1)_j}} \Big[ u^0_{({N_j}-1)_j} - {}^0\bar{C}_{({N_j}-1)_j}\dot{\mathbf{x}}_{2p-3} - {}^0\bar{K}_{({N_j}-1)_j}\mathbf{x}_{2p-3} + {}^0\hat{C}_{{N_j}_j}\dot{\mathbf{x}}_{2p-1} + {}^0\hat{K}_{{N_j}_j}\mathbf{x}_{2p-1} \Big] \\
    \ddot{\mathbf{x}}_{2p} &= 
    \frac{1}{I_{z_{{N_j}_j}}} \Big[ \tau^0_{{N_j}_j} - {}^0\bar{c}_{\theta,{N_j}_j}\dot{\mathbf{x}}_{2p} - {}^0\bar{k}_{\theta,{N_j}_j}\mathbf{x}_{2p} \Big] \notag \\
    &\quad - \frac{1}{I_{z_{({N_j}-1)_j}}} \Big[ \tau^0_{({N_j}-1)_j} - {}^0\bar{c}_{\theta,({N_j}-1)_j}\dot{\mathbf{x}}_{2p-2} - {}^0\bar{k}_{\theta,({N_j}-1)_j}\mathbf{x}_{2p-2} + {}^0\hat{c}_{\theta,{N_j}_j}\dot{\mathbf{x}}_{2p} + {}^0\hat{k}_{\theta,{N_j}_j}\mathbf{x}_{2p} \Big]
\end{align}
